% command to draw a cube, used as illustration in chapter on functions
\newcommand{\drawCube}[4]{
	% first argument: ll position
	% second argument: x shift
	% 3rd argument: # of arguments
	% 4th argument: size (default = 1)
	% define cube coordinates
	\coordinate (A#1) at (#2,    0,     0);
	\coordinate (B#1) at ({#2+#4},0,     0);
	\coordinate (C#1) at ({#2+#4},#4,0);
	\coordinate (D#1) at (#2,    #4,0);
	\coordinate (E#1) at (#2,     0,#4);
	\coordinate (F#1) at ({#2+#4}, 0,#4);
	\coordinate (G#1) at ({#2+#4},#4,#4);
	\coordinate (H#1) at (#2,    #4,#4);

	% draw cube
	\draw[thick] (A#1) -- (B#1) -- (C#1) -- (D#1) -- cycle;
	\draw[thick] (E#1) -- (F#1) -- (G#1) -- (H#1) -- cycle;
	\draw[thick] (A#1) -- (E#1);
	\draw[thick] (B#1) -- (F#1);
	\draw[thick] (C#1) -- (G#1);
	\draw[thick] (D#1) -- (H#1);

	% define lower left corner of doors
	% door with is defined to always be 0.2*size wide and 0.5*size high. We also define that the space between each door is the remaining width,
	% i.e., if we have three doors, then the remaining space is (size-0.2*size*3)/(3+1)

	% calculate how much empty space to leave between each door
	\pgfmathsetmacro\emptyspace{(#4-0.2*#4*#3)/(#3+1)}

	% define amount of space per door
	\pgfmathsetmacro\doorsspace{0.2*#4}

	\foreach \x in {1,...,#3} {
			\pgfmathsetmacro\llCord{\emptyspace*\x+\doorsspace*(\x-1)}
			% lower door coordinate
			\coordinate (llD\x#1) at ({\llCord+#2},0,#4);

			% draw doors
			\fill[blue!50, rounded corners] (llD\x#1) -- ++(0,{#4/2},0) -- ++({0.2*#4},0,0)node[pos=.5, below](in#1\x){\textcolor{blue}{\texttt{x\x}}} -- ++(0,{-#4/2},0);

			% lower door coordinate for value (for later annotation)
			\coordinate (llD\x#1val) at ({\llCord-.1*#4+#2},{-.2*#4},#4); % the coordinate is simple .1*size to the left of the ll corner and .2*size below.
		}

	% coordinate for result door
	\coordinate (llDOut#1) at ({#2+#4},0,{#4*.2});
	\fill[green!50, rounded corners] (llDOut#1) -- ++(0,{#4/2},0) -- ++(0,0,{#4/2})node[pos=0.8,rotate=45, below, align=center](out#1){\textcolor{black}{\tiny\texttt{out#1}\normalsize}} -- ++(0,{-#4/2},0);

	% draw a 2D rectangle around the cube
	\node[fit=(E#1) (C#1), inner sep=4pt, red, dotted, draw, line width=1pt,rounded corners](dottedRect#1) {};
	% annotate the two-d rect
	\node[above=5pt of dottedRect#1.north west, anchor=west, align=left]{\textcolor{red}{\texttt{def a#1(\textcolor{blue}{\foreach\x in {1,...,#3}{x\x, }...})}\normalsize}};
}

\newgeometry{left=1.5cm,right=1.5cm,top=2cm,bottom=2.5cm}
\fancyfootoffset{0pt}% <- recalculate \headwidth

\section{\thesubject: \thetopic\\\normalsize
\mainauthor}

\subsection{Alma-Bereich 1: Lernende --- Fachwissenschaftliche und FD-Kompetenzen: Lerninhalte adäquat aufbereiten}

\textbf{Fragestellung 1}: Wie kann es gelingen, komplexe fachliche Inhalte auf solche Weise zu vermitteln, dass diese für die SuS verständlich werden?

Diese Fragestellung beschäftigt mich seit Beginn meiner Unterrichtstätigkeiten vor eineinhalb Jahren. Die nachfolgende Textbox fasst eine Situation zusammen, welche ich dabei selber erlebt habe.

\begin{myExBox}[title=Beispiel aus der eigenen Unterrichtserfahrung]
Frustriert trat ich ins Zimmer der Prorektorin. Es ging um die Nachbesprechung meiner Unterrichtssequenz, welche sie soeben hospitiert hatte. Obschon ich mir sehr viel Mühe gegeben hatte, war der Inhalt, welcher ich den SuS vermitteln wollte, leider just in dieser Lektion überhaupt nicht angekommen. Meine Prorektorin hospitierte zuvor meine Lektion an einem Donnerstagnachmittag, in der zweiten Lektion einer Doppellektion, direkt nachdem die SuS eine Informatik-Prüfung bei mir in der Vorlektion hatten. Leider waren die SuS (die am Donnerstag auch sonst 10 Lektionen haben und müde sind) fixfertig und konnten die Inhalte kaum aufnehmen.

%Die wohl zu schwer geratene Prüfung belastete nicht nur die SuS, sondern auch mich, da ich sie ja für Ihre Lern-Aufwände gerecht behandeln wollte. Nun kam jedoch auch noch hinzu, dass die Prorektorin im Raum sass, was dazu führte, dass ich in eine Art ``Notmodus'' schaltete und meine umfangreichen Vorbereitungen demonstrieren und schnell zum Inhalt der Lektion kommen wollte.

Für die Vorbereitung der neuen Inhalte hatte ich mir sehr viel Mühe gegeben: Mein Referat zum Thema Funktionen beinhaltete Unterlagen, für deren Einarbeitung ich mehrere Tage gebraucht hatte. Schon nur für die Grafiken hatte ich mehrere Nachmittage investiert, etwa folgende Grafik, die Funktionen bildlich mit Abteilungen in einer Fabrik vergleicht:

\begin{figure}[H]
	% Use the \drawCube macro to draw various cubes
	\centering
	\begin{tikzpicture}
		\drawCube{1}{0}{2}{3}
		% feed input values into a1
		\node[fill=blue!50, rounded corners] (input1) at (llD11val){3};
		\node[fill=blue!50, rounded corners] (input2) at (llD21val){12};
		\draw[thick,->] (input1) -- (in11);
		\draw[thick,->] (input2) -- (in12);
		\node[fill=green!50, rounded corners](outVal1) at (5,-.6,3){36};
		\draw[thick,->] (out1) -- (outVal1);
		\drawCube{2}{5.25}{3}{3}
		\drawCube{3}{10.5}{2}{3}
		\draw[thick,->,rounded corners] (out1) -- ++(0,-3.25,0) -- ++(3.15,0,0)node[below]{\lstinline{return out1}} -- ++(4.65,0,0) -- (in31);
	\end{tikzpicture}
	\caption{Illustration einer Code-Struktur, bei welcher mehrere Funktionen zusammenarbeiten}
\end{figure}
In der Hitze des Gefechts versuchte ich zwar, einige Beispiele und Analogien von Funktionen im echten Leben aufzuzeigen (Modularität in einer Fabrik, Kochrezepte etc.), allerdings waren es wohl etwas viele Beispiele, zudem wurden die SuS leider zu wenig in den Prozess involviert. Meine Vorgehensweise, trotz grossem Engagement, war wohl nicht zielführend. So kam es dazu, dass sich während der Lektion herausstellte, dass ein Grossteil der SuS meinen Erklärungen leider nicht gefolgt hatte und nochmals persönlich bei mir nachfragen musste. Die Übungen wurden grösstenteils nicht gelöst. Dies fiel auch der Prorektorin auf und führte meinerseits zu grosser Frustration, insbesondere angesichts meiner investierten zeitlichen Ressourcen.

\end{myExBox}

Im Rahmen der Analysen dieser Situation ergaben sich folgende Erkenntnisse: Auch eine aufwendig und sorgfältig vorbereitete Lektion kann bei den Lernenden ins Leere laufen, wenn deren Belastungsgrenze bereits überschritten ist. Die intensive kognitive Beanspruchung der SuS nach einer Prüfung und am Ende eines langen Schultages führte dazu, dass zentrale Inhalte nicht aufgenommen werden konnten – trotz visualisierten Materialien und durchdachter Strukturierung. Die Reflexion der Situation machte deutlich, dass fachlich gut aufbereitete Inhalte allein nicht genügen, wenn sie nicht auf die aktuelle Verfassung und Motivation der Lernenden abgestimmt sind.

Zentral ist dabei die Bedeutung der Motivation für gelingende Lernprozesse. Gemäss \textcite{ryan, deci} ist Motivation nicht bloss ein Zustand, sondern ein Ergebnis der Passung zwischen Person und Situation \cite{ryan, deci, moeller}. Wenn das Bedürfnis nach Kompetenz, Autonomie und sozialer Eingebundenheit nicht erfüllt ist, entsteht keine intrinsische Motivation. In der beschriebenen Situation war vor allem das Kompetenzempfinden beeinträchtigt – sowohl bei den SuS als auch bei mir als Lehrperson. Die Enttäuschung darüber, dass eine gut vorbereitete Lektion \enquote{nicht funktioniert}, war gross.

Für meine eigene Entwicklung als Lehrperson nehme ich aus dieser Situation mit, dass Lerninhalte nicht nur fachlich korrekt und strukturiert aufbereitet, sondern stets auch in Bezug zu den situativen und motivationalen Voraussetzungen der Lernenden gesetzt werden müssen. Es braucht ein feines Gespür dafür, wann Reduktion, Wiederholung oder auch eine Änderung der Zielsetzung angebracht sind. Wie Jenni und Thyen betonen, geht es im professionellen Handeln von Lehrpersonen nicht nur um das Reagieren auf Störungen, sondern um das Erkennen und Verstehen von Situationen im Kontext pädagogischer Beziehungen \cite{jenni,thyen}. In diesem Sinne wird die Verbindung von fachlicher und fachdidaktischer Kompetenz mit pädagogischer Sensibilität zu einem Schlüssel für nachhaltiges Lernen.

Diese Erkenntnisse tragen im Sinne des Alma-Bereichs 1 dazu bei, meine Fähigkeit zu stärken, fachwissenschaftliche Inhalte lernwirksam und situationsgerecht aufzubereiten.


\subsection{Alma-Bereich 2: Lehrperson --- ESW- und FD-Kompetenzen: Lehrhandeln ethisch reflektieren}
\textbf{Fragestellung 2}: Wie gelingt es mir als Lehrperson in einem staatlich-öffentlichen Umfeld, regelmässiges und ehrliches Feedback zu mir als Lehrperson zu erhalten, um mich weiterzuentwickeln?

Lehrpersonen sind gemäss Hattie (2009) dann besonders wirksam, wenn sie systematisch Feedback einholen und daraus gezielte Anpassungen ableiten \cite{alma991171885796805501}. Ein zentrales Anliegen in meiner bisherigen Entwicklung als Lehrperson ist es, besser zu verstehen, wie mein Unterricht bei den Lernenden ankommt und wie ich als Person im Klassenzimmer wahrgenommen werde. Trotz bereits unternommener Versuche, die Hemmschwelle für Rückmeldungen zu senken, erhalte ich für meinen Geschmack noch immer zu wenig differenziertes und konstruktives Feedback. Gerade weil mir die kontinuierliche persönliche und professionelle Weiterentwicklung ein grosses Anliegen sind, beschäftigt mich die Frage, wie ich an die bereits bestehenden Gedanken und Eindrücke der Lernenden über mich als Lehrperson herankomme – weniger aus einer Haltung der Unsicherheit heraus, sondern vielmehr aus einem echten Interesse an Entwicklung und Selbstreflexion.

Die bisherigen Rückmeldungen – etwa aus anonym durchgeführten Umfragen gegen Ende einer Unterrichtseinheit oder am Ende eines Praktikums – haben mir gezeigt, wie wertvoll gezieltes Feedback sein kann. Besonders hilfreich war eine einfache, offene Frage: \enquote{Was kann ich besser machen?} (auch wenn ich diese manuell auswerten musste), Die Rückmeldungen darauf waren ehrlich, überraschend und haben mir neue Perspektiven auf mein Handeln eröffnet. Gleichzeitig wurde mir bewusst, dass es nicht zwingend ein aufwendiges Auswertungssystem braucht, sondern vielmehr persönliche Nähe und ehrlicher Austausch mit den SuS.

Ein zentrales Lernfeld besteht für mich darin, die Feedbackerhebung stärker zu strukturieren und in den Unterrichtsverlauf zu integrieren. Statt Rückmeldungen erst am Ende eines Semesters einzuholen, plane ich, diese künftig unmittelbar nach abgeschlossenen Unterrichtseinheiten oder thematischen Blöcken durchzuführen – also dann, wenn die Eindrücke noch frisch sind. Denkbar wäre auch, Feedback thematisch einzubetten, z.\,B. im Informatikunterricht durch eine spielerische Bewertung der erarbeiteten Produkte oder einer Peer-Review-Phase, bei der auch die Lehrperson einbezogen wird. Die Erkenntnis, dass Feedback nicht immer im Rahmen einer Plenumsdiskussion stattfinden muss – da diese oft wenig ergiebig ist –, hat sich für mich als hilfreich erwiesen.

Ich nehme mit, dass der Weg zu mehr Rückmeldung von einem klaren Rahmen und einer gewissen Regelmässigkeit begleitet sein sollte. Hierzu braucht es weiterhin meine Offenheit -- nicht nur für Lob, sondern auch für kritische Hinweise. Die Sorge vor einer möglichen Stagnation ist für mich ein Antrieb, mein professionelles Handeln kontinuierlich zu hinterfragen. Rückmeldungen der SuS bieten dabei eine wertvolle Ressource, um blinde Flecken zu erkennen und meine Unterrichtsgestaltung gezielt zu verbessern.

\subsection{Alma-Bereich 3: Lernangebote --- Transversale Kompetenzen: Digitalität integrieren}
\textbf{Fragestellung 3}: Wie kann ich digitale Inhalte so einsetzen, dass sie den (analogen) Lernprozess begünstigen, bzw. erweitern?

Als Informatik-Lehrperson liegt es aufgrund meines Fachs nahe, dass ich viele Inhalte mit digitalen Mitteln vermittle. Dabei ist es mir jedoch stets ein Anliegen, diese so zu wählen, dass sie die Lernprozesse in meinen Unterrichtseinheiten sinnvoll unterstützen. Häufig verwendete digitale Lernmittel in meinem Unterricht sind unten kurz aufgelistet und erklärt.

\textbf{Moodle }ist eine weitverbreitete, modular aufgebaute Lernplattform (Learning Management System, LMS), die es Lehrpersonen ermöglicht, digitale Lernumgebungen strukturiert und vielseitig zu gestalten. Im Informatikunterricht bietet Moodle besonders viele Vorteile, da es nicht nur als Ablage- und Kommunikationsplattform dient, sondern auch interaktive Elemente und automatisierte Auswertungen unterstützt. Die Plattform fördert dabei sowohl selbstgesteuertes Lernen als auch transparente Kommunikation über Lernziele und Leistungsanforderungen. 

Besonders nützlich für den Informatikunterricht sind folgende Funktionen:

\begin{itemize}
  \item \textbf{CodeRunner-Plugin:} Mit diesem Plugin können automatisch bewertbare Programmieraufgaben direkt in Moodle eingebunden werden. Lernende erhalten unmittelbar Feedback zu ihrer Lösung, was iterative Lernprozesse unterstützt.
  \item \textbf{Abgabe von Aufgaben:} Schülerinnen und Schüler können Programmcode, Dokumentationen oder Projektdateien direkt hochladen. Lehrpersonen haben die Möglichkeit, strukturierte Rückmeldungen zu geben und Korrekturen effizient zu verwalten.
  \item \textbf{Unterlagen bereitstellen:} Materialien wie Theorieblätter, interaktive Übungen oder Screencasts lassen sich geordnet zur Verfügung stellen, was den Zugang zu Lerninhalten erleichtert und eine strukturierte Nachbearbeitung ermöglicht.
  \item \textbf{Foren zur Diskussion:} Über Foren können Lernende Fragen stellen, sich gegenseitig unterstützen oder vertiefende Diskussionen führen. Dies fördert kollaboratives Lernen und entlastet die Lehrperson bei häufig gestellten Fragen.
  \item \textbf{Online-Prüfungen mit Safe Exam Browser (SEB):} Moodle lässt sich in Kombination mit dem Safe Exam Browser für sichere, digital durchgeführte Prüfungen nutzen. Dies ermöglicht eine flexible Leistungsüberprüfung mit erhöhtem Schutz vor Täuschung.
\end{itemize}

Durch die Kombination dieser Werkzeuge bietet Moodle eine lernförderliche Umgebung, die sowohl individuelle als auch kooperative Lernprozesse unterstützt. Lehrpersonen können Lernfortschritte gezielter begleiten, während Lernende mehr Verantwortung für ihren Lernprozess übernehmen können. Insbesondere im Informatikunterricht, wo digitale Kompetenzen zentral sind.

Neben den Moodle-eigenen Funktionen setze ich im Informatikunterricht eine Reihe zusätzlicher digitaler Werkzeuge ein, die sowohl der Veranschaulichung als auch der Interaktivität dienen:

\begin{itemize}
  \item \textbf{H5P-Snippets:} Diese ermöglichen die Integration interaktiver Elemente wie Multiple-Choice-Aufgaben, Drag-and-Drop oder interaktive Videos direkt in Moodle. Sie sind besonders hilfreich für formative Überprüfungen und die Selbstkontrolle.
  \item \textbf{Interaktive Webseiten für Kryptologie:} Für das Thema Kryptologie verwende ich beispielsweise die Webseite \href{https://cryptbreaker.marcwidmer.xyz}{cryptbreaker.marcwidmer.xyz}, auf der Lernende spielerisch Verschlüsselungsverfahren analysieren und knacken können. Dies fördert das entdeckende Lernen und die Anwendung theoretischer Konzepte in einer realitätsnahen Umgebung.
  \item \textbf{LaTeX für Skripts und Präsentationen:} Für die Erstellung von Unterrichtsskripten und Folien nutze ich LaTeX, da es eine präzise und gut strukturierte Darstellung komplexer Inhalte – insbesondere im Bereich der formalen Sprachen und der Mathematik – erlaubt. 
  \item \textbf{Künstliche Intelligenz (z.\,B. ChatGPT):} Tools wie ChatGPT helfen mir bei der schnellen Formulierung von Erklärtexten, beim Generieren von Beispielcode oder beim Erstellen didaktischer Materialien. Die Kombination von AI-Unterstützung und strukturiertem Layout (z.\,B. in LaTeX) ermöglicht eine erhebliche Effizienzsteigerung im Vorbereitungsalltag.
\end{itemize}

Didaktisch betrachtet zeigen diese digitalen Werkzeuge grosses Potenzial – nicht primär, weil sie den Lernenden den Stoff besser erklären, sondern weil sie mir als Lehrperson helfen, meinen Unterricht flexibler, individueller und zeiteffizienter zu gestalten. Die zentrale Rolle der Lehrperson bleibt dabei unbestritten: Viele Schülerinnen und Schüler scheitern nicht an der Komplexität der Inhalte, sondern bereits an vermeintlich einfachen Hürden wie Syntaxfehlern oder der Interpretation einer Aufgabenstellung. Digitale Umgebungen wie CodeRunner oder H5P können hier sofortiges, gezieltes Feedback geben – insbesondere bei formalen Fehlern oder grundlegenden Verständnisproblemen.

Dennoch ersetzt kein Tool die professionelle Begleitung, das gezielte Nachfragen oder die situationsgerechte Erklärung durch eine Lehrperson. Vielmehr verschaffen mir diese digitalen Helfer die nötigen Freiräume, um mich stärker auf individuelle Unterstützung zu konzentrieren. In diesem Sinne sind digitale Tools in erster Linie ein \enquote{game changer} für uns Lehrpersonen – sie vereinfachen unsere Arbeit, strukturieren den Unterricht und unterstützen die Differenzierung. Für die Lernenden stellen sie hilfreiche Ergänzungen dar, ersetzen jedoch nicht die Notwendigkeit einer klaren, verständnisorientierten Vermittlung durch eine \enquote{real} präsente ansprechbare Lehrperson.

Im Idealfall tragen diese digitalen Werkzeuge jedoch auch zur Förderung zentraler pädagogischer Ziele bei, wie der Stärkung von \emph{Selbstwirksamkeit}, \emph{Selbstbestimmtheit} und \emph{Selbstkompetenz} bei den Lernenden. Indem sie individuelle Rückmeldung erhalten, sich Inhalte im eigenen Tempo erarbeiten können und zunehmend Verantwortung für ihren Lernprozess übernehmen, werden diese drei Dimensionen unterstützt. Damit greifen die digitalen Elemente nicht nur in den methodischen Alltag ein, sondern leisten auch einen Beitrag zur Persönlichkeitsentwicklung \cite{alma991170251641605501}.



\appendix
\section{Anhang}
\subsection{KI-Deklaration}

Der gesamte Text basiert auf eigenen Erfahrungen und Überlegungen. Künstliche Intelligenz in Form von Sprachmodellen wurde verwendet, um die konzise Formulierung des Text zu unterstützen. Alle Äusserungen in diesem Text wurden persönlich auf ihre inhaltliche Richtigkeit und stilistische Kohärenz und Authentizität überprüft. Alle gemachten Äusserungen entsprechen daher meinen eigenen Erfahrungen. Dieser Abschnitt wurde gänzlich ohne künstliche Intelligenz verfasst.

\printbibliography[heading=bibsubsection]